\scanpage{265}
\mspage{6}

\deleted{le système de signes \uline{opposé} de celui de $I$. \inserted{\unsure{Je vais expliciter}}}

NB Si $i$ n'est pas un bout, et si néanmoins l'une des deux catégories $i\backslash I$, $I_{/i}$ est ponctuelle, disons $I_{/i}$, alors au voisinage immédiat de $i$ $I$ a la forme

\diagram{\begin{tikzcd}[ampersand replacement=\&,column sep=large]
\underset{i-1}{\vphantom{\bullet}} \& \underset{i}{\bullet}\arrow[l,"{\underline{u}_{\deleted{\unsure{i-1|i+1}}\unsure{i}}}"']\arrow[r,"{\underline{u}_{i+1}}"] \& \underset{i+1}{\vphantom{\bullet}}
\end{tikzcd}}

\deleted{(\unsure{$u_i$ et $u_{i+1}$}} ($u_i$ rétrograde, $u_{i+1}$ directe -- si c'était $i\backslash I$ la ponctuelle, alors $u_i$ serait directe, $u_{i+1}$ rétrograde -- en tous les cas \deleted{\unsure{deux}} \unsure{ces} deux flèches sont de sens opposé, i.e. non composables). \deleted{\unsure{alors}} On trouve donc que $i\backslash I$ est une catégorie de la forme

\diagram{\begin{tikzcd}[ampersand replacement=\&,row sep=small,column sep=small]
 \& \& i\arrow[dl]\arrow[dr] \& \& \\
 \& i-1\arrow[ddl,dotted,no head]\arrow[ddl,shorten <=1.4em] \& \& i+1\arrow[dr,dotted,no head]\arrow[dr,shorten <=1.1em] \& \\
 \& \& \& \& i+\unsure{s} \\
 i-r \& \& \& \&
\end{tikzcd}}

où $r$ est la longueur \deleted{\unsure{entière}} de la chaîne structurale contenant $\underline{u}_i$, $s$ celle de la chaîne structurale contenant $\underline{u}_{i+1}$. Cette catégorie est une somme amalgamée $\mathbb{\Delta}_r\amalg\mathbb{\Delta}_s$ (recollées suivant leur origine), mais n'est pas isom. à un $\mathbb{\Delta}_j$, p.ex. faute d'avoir un objet final.

Cette observation prouve a). Pour prouver b, il suffit de prouver que l'ordre total sur $I$ est déterminé de façon intrinsèque, \deleted{\unsure{On va voir}} \deleted{\unsure{Pour|Donc} ceci} par \unsure{le|les} choix de l'\uline{origine} \inserted{$s_0$} de $I$ (parmi ces deux extrémités). Supposons que l'on ait déjà déterminé intrinsèquement les sommets $s_0,s_1,\ldots s_i$ ($i\unsure{\geq|>}0$). Comment déterminer $s_{i+1}$ ? \deleted{(On regarde \unsure{$s_i$})} \inserted{(quand $i<n$)} Soit \deleted{$J=\{s_0,\ldots s_{i-1}\}$} $I'$ la sous-catégorie pleine de $I$

% archive sheet number: 264
